\pdfvariable trailerid {[<C2922026090200000000000000000000><C2922026090200000000000000000000>]}
\documentclass[10pt]{article}
\usepackage[margin=0.78in]{geometry}
\usepackage{amsmath,amssymb,amsthm,booktabs,microtype,hyperref}
\hypersetup{colorlinks=true,linkcolor=blue,urlcolor=blue}
\providecommand{\CRevisionRound}{2}
\newtheorem{theorem}{Theorem}
\newtheorem{lemma}{Lemma}
\newtheorem{proposition}{Proposition}
\newcommand{\R}{\mathbb R}
\newcommand{\D}{\mathrm d}
\newcommand{\one}{\mathbf 1}
\title{All-Event Dynamics of Finite Sticky Particles:\\
Weighted Isotonic Geometry, Dissipation, and Weak Closure}
\author{Route-A source-local certificate HCS-C292}
\date{2 September 2026}
\begin{document}
\maketitle

\begin{abstract}
We close the forward dynamics of an arbitrary finite one-dimensional system
of positive point masses undergoing perfectly inelastic sticking.
Initially coincident labels are canonically premerged, after which a single
construction handles binary, multi-cluster, and spatially disjoint
simultaneous events.
\ifnum\CRevisionRound=0
The position vector is the mass-weighted isotonic projection of the free
configuration, which proves uniqueness and the no-splitting property.
\fi
\ifnum\CRevisionRound>0
Equivalently, positions are slopes of a lower convex hull in cumulative-mass
coordinates.  Every event conserves mass and momentum and has an exact
pairwise-variance kinetic-energy loss.
\fi
\ifnum\CRevisionRound>1
The induced atomic density and momentum solve pressureless Euler in the
distributional sense and obey the kinetic-energy entropy inequality.
Strict independent evidence exercises all event types without substituting
finite computation for the all-data proof.
\fi
\end{abstract}

\section{Model, conventions, and source ownership}
Let raw labels carry masses $m_i>0$, ordered positions
$x_1\leq\cdots\leq x_N$, and velocities $v_i\in\R$.  All labels at the same
initial position are replaced at time zero by one canonical particle whose
mass and velocity are
\begin{equation}\label{eq:premerge}
 M=\sum_i m_i,
 \qquad V=M^{-1}\sum_i m_i v_i.
\end{equation}
The resulting canonical positions are strictly increasing.  This convention
preserves mass and momentum and regards any raw-label energy drop as part of
initialization.  Zero mass is excluded because both~\eqref{eq:premerge} and
the weighted geometry would otherwise have undefined blocks.

Canonical particles move ballistically until contact.  Every maximal
consecutive block occupying one position then becomes one particle, with
mass and velocity given by~\eqref{eq:premerge}; it never splits.  Velocities
at event times are right-continuous.  This is a forward entropy dynamics:
we make no backward-uniqueness claim.

The monotone-cone formulation and sticky semigroup are classical; a direct
owner is Natile and Savar\'e~\cite{NS2009}.  Brenier, Gangbo, Savar\'e, and
Westdickenberg treat sticky Lagrangian pressureless flows with interactions
\cite{BGSW2013}, and Hynd gives a Lagrangian pressureless-Euler construction
\cite{Hynd2019}.  The result below is a convention-complete, independently
audited workspace certificate, not a claim of literature priority.

\section{The all-event theorem}
Write the canonical data again as $(m_i,x_i,v_i)_{i=1}^n$, now with
$x_1<\cdots<x_n$.  Give $\R^n$ the weighted norm
$\lVert z\rVert_m^2=\sum_i m_i z_i^2$ and set
\[
 K=\{z\in\R^n:z_1\leq\cdots\leq z_n\}.
\]
Let $P_K^m$ denote the unique projection in this norm.

\begin{theorem}[Arbitrary finite all-event flow]\label{thm:main}
For every canonical finite datum there is exactly one global forward sticky
flow.  Its labelled positions are
\begin{equation}\label{eq:projection}
 X(t)=P_K^m(x+t v),\qquad t\geq0.
\end{equation}
The equal-coordinate blocks of $X(t)$ are precisely the physical clusters;
their velocity is the mass-weighted average of the $v_i$ in that block.
The partitions only coarsen.  Thus there are at most $n-1$ nontrivial
mergers, while any one event time may contain several disjoint contact
positions and any contact position may contain more than two incoming
clusters.
\end{theorem}

\begin{proof}
The finite event construction is explicit.  At any state, compare adjacent
clusters $A,B$.  If $V_A>V_B$, their prospective delay is
\[
 \tau_{A,B}=\frac{X_B-X_A}{V_A-V_B};
\]
otherwise it is infinite.  Advance every cluster through the least finite
delay.  At that time merge \emph{all} maximal consecutive collocated blocks,
including blocks at different positions.  Each nontrivial merge lowers the
cluster count, so the process is global after at most $n-1$ mergers.  This
also covers several adjacent prospective delays becoming equal.

It remains to identify the construction and prove uniqueness.  Weighted
pool-adjacent-violators starts from the free values $y_i(t)=x_i+t v_i$ and
replaces every adjacent inversion by its mass-weighted mean.  Its terminal
blocks satisfy the variational inequalities for projection onto $K$, hence
give the unique vector~\eqref{eq:projection}.  Before the first contact no
pool is needed.  At first contact, exactly the maximal collocated incoming
blocks are pooled, and their slope in $t$ is their barycentric velocity.
Repeating this argument after the event identifies the projection with the
event construction on the next interval.  Induction closes every interval.
Once a set of indices is pooled, the same barycentric rule governs it in the
next interval, so a physical block cannot split.  The projection therefore
selects the unique forward sticky continuation.
\end{proof}

Separated clusters with equal velocity have infinite prospective delay and
do not collide.  A collision may fail to occur altogether.  Simultaneity is
not resolved by arbitrary pair ordering: maximal contact blocks are merged
in one operation.

\ifnum\CRevisionRound>0
\section{Lower convex hull and event geometry}
Put $M_0=S_0=0$ and, for $1\leq j\leq n$,
\begin{equation}\label{eq:cumsum}
 M_j=\sum_{i\leq j}m_i,
 \qquad S_j(t)=\sum_{i\leq j}m_i(x_i+t v_i).
\end{equation}
Take the greatest convex minorant of the polygonal line through
$(M_j,S_j(t))_{j=0}^n$.

\begin{proposition}[Hull representation]\label{prop:hull}
On each interval $[M_{a-1},M_b]$ where the minorant has one affine face, its
slope equals $X_i(t)$ for every $a\leq i\leq b$.  Its faces are therefore
exactly the sticky clusters at time $t$.
\end{proposition}
\begin{proof}
The slope of the chord from $M_{a-1}$ to $M_b$ is the weighted mean of
$y_a(t),\ldots,y_b(t)$.  Convexity orders successive slopes.  Maximal faces
therefore give a feasible isotonic vector.  The minorant inequalities are
the prefix-sum form of the projection variational inequality; equality at
face endpoints supplies complementary slackness.  Strict convexity of the
weighted square norm identifies this vector with~\eqref{eq:projection}.
\end{proof}

The hull picture also makes simultaneous events unambiguous.  At an event,
one or several neighboring groups of faces become collinear.  Replacing each
group by its maximal common face produces the unique post-event partition.
Because later physical partitions coarsen, an equality face is never split
by a label-level tie convention.

\section{Exact conservation and dissipation}
Suppose one event group contains incoming clusters with masses $M_a>0$ and
velocities $V_a$, $1\leq a\leq r$.  Set
$M=\sum_aM_a$, $P=\sum_aM_aV_a$, and $\bar V=P/M$.

\begin{proposition}[Multi-cluster loss identity]\label{prop:loss}
Mass and momentum are conserved, and the kinetic-energy loss is
\begin{equation}\label{eq:loss}
 \Delta E
 =\frac12\sum_aM_aV_a^2-\frac12M\bar V^2
 =\frac1{2M}\sum_{a<b}M_aM_b(V_a-V_b)^2\geq0.
\end{equation}
The formula adds over disjoint contact positions at the same event time.
Globally, total mass and momentum are constant, and
\begin{equation}\label{eq:com}
 \sum_C M_CX_C(t)=\sum_i m_ix_i+t\sum_i m_iv_i.
\end{equation}
\end{proposition}
\begin{proof}
The definitions of $M,P$, and $\bar V$ prove mass and momentum conservation.
Expanding
$\sum_aM_a(V_a-\bar V)^2$ gives the first difference
in~\eqref{eq:loss}; symmetrizing gives the pairwise expression.  Each term is
nonnegative.  Ballistic differentiation proves~\eqref{eq:com} between
events, while event continuity and momentum conservation remove every jump.
Disjoint event groups have disjoint incoming index sets, so their losses add.
\end{proof}

This exact identity includes binary and arbitrary simultaneous mergers.  It
also distinguishes the physical entropy selection from a merely
mass-conserving coalescence rule.
\fi

\ifnum\CRevisionRound>1
\section{Atomic pressureless Euler closure}
At non-event times let the sum run over current clusters and define Radon
measures
\begin{equation}\label{eq:measures}
 \rho_t=\sum_C M_C\delta_{X_C(t)},\quad
 j_t=\sum_C M_CV_C\delta_{X_C(t)},\quad
 q_t=\sum_C M_CV_C^2\delta_{X_C(t)}.
\end{equation}
The value assigned to $V_C$ at finitely many event instants is immaterial in
space--time integration; we retain the right-continuous convention.

\begin{theorem}[Weak pressureless Euler and entropy]\label{thm:pde}
For all compactly supported smooth test functions $\phi,\psi$ on
$[0,\infty)\times\R$,
\begin{align}
 \int_0^\infty\!\!\int
   (\partial_t\phi\,\D\rho_t+\partial_x\phi\,\D j_t)\D t
   +\int\phi(0,x)\,\D\rho_0(x)&=0,\label{eq:massweak}\\
 \int_0^\infty\!\!\int
   (\partial_t\psi\,\D j_t+\partial_x\psi\,\D q_t)\D t
   +\int\psi(0,x)\,\D j_0(x)&=0.\label{eq:momweak}
\end{align}
Moreover, with
$e_t=\frac12\sum_CM_CV_C^2\delta_{X_C(t)}$ and
$f_t=\frac12\sum_CM_CV_C^3\delta_{X_C(t)}$,
\begin{equation}\label{eq:entropy}
 \partial_t e+\partial_x f
 =-\sum_{\tau}\sum_{g\text{ at }\tau}
   \Delta E_{\tau,g}\,\delta_{(\tau,X_{\tau,g})}\leq0
\end{equation}
as a distribution.
\end{theorem}
\begin{proof}
On every open ballistic segment,
$\frac{\D}{\D t}\phi(t,X_C(t))=\phi_t+V_C\phi_x$.
Multiplying by $M_C$, integrating segment by segment, and summing gives
\eqref{eq:massweak}.  Endpoint atoms at a collision cancel because incoming
and outgoing masses agree.  The same calculation with weights $M_CV_C$
gives~\eqref{eq:momweak}; its endpoint atoms cancel by momentum conservation.
Using weights $M_CV_C^2/2$ instead leaves at each event exactly the negative
of~\eqref{eq:loss}, proving~\eqref{eq:entropy}.
\end{proof}

Thus the event construction is not only a particle algorithm: it closes both
conservation laws and the entropy defect for arbitrary finite data.

\section{Executable evidence and boundary atlas}
The deterministic receipt contains six scenarios: a two-stage binary
merger, a simultaneous three-cluster merger, two disjoint simultaneous
mergers, initial coincident labels, a no-collision face, and a cascade.  It
records 20 raw labels, 18 canonical particles, 18 premerge cells, seven
event times, eight event groups, 75 projection cells, six conservation
cells, and eight weak-balance cells.  An independent checker enumerates all
contiguous partitions rather than importing the event producer; it reports
1,538 assertions.  SymPy checks 255 exact identities, two isolated evidence
replays are byte identical, and 66/66 mutations are rejected: 45 evidence-JSON
and 21 evaluation-YAML attacks.  The latter include a duplicate-key-rejecting
evaluation YAML semantic lock.  These are regression receipts, not a finite
proof of Theorems~\ref{thm:main}--\ref{thm:pde}.

The boundaries are explicit: masses are strictly positive; coincident
initial labels are premerged; separated equal-speed clusters need not meet;
multiple positions may collide at one time; a contact may have arbitrarily
many incoming clusters; only forward uniqueness is selected; and the finite
system has no accumulation of collision times because every merger reduces
the cluster count.

\section{Route-A assessment and disclosures}
All five Route-A axes fail, so the strict tuple is
\[
 \texttt{(A0\_FAIL,A1\_FAIL,A2\_FAIL,A3\_FAIL,A4\_FAIL)}.
\]
The verdict is \texttt{ROUTE\_A\_REJECTED}; obstruction \texttt{HEN-O276}
is closed only as a source-dynamics theorem, and Route B is disabled.  The
literal scope is \texttt{NO\_BAD\_EULER\_OR\_ROOT\_NUMBER}.  No arithmetic
local data, Euler factors, root numbers, automorphy, target divisor law,
functional equation, zero match, or Hilbert--P\'olya operator is claimed.

All source, evidence, checker, mutation, and build instructions are included
in the package.  No personal or sensitive data and no human or animal
subjects are involved; no ethics approval is required.  The author reports
no conflict of interest and no external funding.  CRediT roles are
conceptualization, formal analysis, software, validation, visualization,
writing, and reproducibility curation.  AI assistance was used for drafting,
symbolic code, and consistency checks; the mathematical claims were checked
against the displayed proofs and independently executable gates.
\fi

\begin{thebibliography}{3}
\bibitem{NS2009}
L. Natile and G. Savar\'e,
``A Wasserstein approach to the one-dimensional sticky particle system,''
\emph{SIAM J. Math. Anal.} 41 (2009), 1340--1365,
\href{https://arxiv.org/abs/0902.4373}{arXiv:0902.4373}.
\bibitem{BGSW2013}
Y. Brenier, W. Gangbo, G. Savar\'e, and M. Westdickenberg,
``Sticky particle dynamics with interactions,''
\emph{J. Math. Pures Appl.} 99 (2013), 577--617,
\href{https://arxiv.org/abs/1201.2350}{arXiv:1201.2350}.
\bibitem{Hynd2019}
R. Hynd, ``Lagrangian Coordinates for the Sticky Particle System,''
\emph{SIAM J. Math. Anal.} 51 (2019), 3769--3795,
\href{https://doi.org/10.1137/19M1241775}{doi:10.1137/19M1241775}.
\end{thebibliography}
\end{document}
